\section{\termDeformationOfFloerEquations{}}
\label{sec:deformation-of-floer-equations}

迄今我们考察的是哈密顿扰动的\personCauchy{}--\personRiemann{}方程~\eqref{eq:floer-equation-zero-conormal}；按照 (Oh97b)，我们称它为\emph{\termDynamicalVersion{}}。

另一方面，也可以考虑\emph{真正的}\personCauchy{}--\personRiemann{}方程
\begin{equation}
  \begin{cases}
    \dfrac{\partial v}{\partial\tau}
    +J^H\dfrac{\partial v}{\partial t}=0,\\[6pt]
    v(\tau,0)\in\phi_H^1(o_N),
    \qquad v(\tau,1)\in\nu^*S,
  \end{cases}
\label{eq:geometric-floer-equation}
\end{equation}
这里把 $v$ 看作道路
\[
  v:\mathbb R\longrightarrow\Omega(L,\nu^*S),
\]
其中
\[
  L=\phi_H^1(o_N)
\]
且
\[
  \Omega(L,\nu^*S)
  =\{\gamma:[0,1]\to T^*N
      \mid \gamma(0)\in L,\ \gamma(1)\in\nu^*S\},
\]
而
\[
  J_t^H
  =\bigl(\phi_H^t(\phi_H^1)^{-1}\bigr)_*J_t.
\]
我们称这一版本为\emph{\termGeometricVersion{}}。\footnote{译者注：原书公式 (12.7.36) 使用未知映射 $v$，紧随其后的说明却称道路为 $u$；译文统一为 $v$。}

现在更详细地描述\personFloer{}同调的几何版本。要点在于，拉格朗日\personFloer{}理论的动力学版本与几何版本之间存在\termFiltrationPreservingIsomorphism{}。

把有限能量解的集合记为
\[
  \widetilde{\mathcal M}(L_H,\nu^*S;J^H),
\]
并把它关于 $\mathbb R$-平移的商记为
\[
  \mathcal M(L_H,\nu^*S;J^H).
\]
由此得到\personFloer{}同调的几何版本
\[
  HF_*\bigl(\phi_H^1(o_N),\nu^*S;J\bigr),
\]
其生成元是交点
\[
  \phi_H^1(o_N)\cap\nu^*S.
\]
这一版本的优点，是它只依赖于拉格朗日子流形 $L=\phi_H^1(o_N)$，因而对 $H$ 的依赖很弱。\footnote{译者注：本段的方程与模空间使用 $J^H$，原书在相伴\personFloer{}同调中却改记 $J$；译文照录，不自行判定此处是否应为 $J^H$。}

下述引理不难直接验证，但十分关键。

\begin{lemma}
\label{lem:dynamical-geometric-floer-correspondence}
设\footnote{译者注：原书把这里的零截面写成此前未定义的 $o_S$；本节前文与下述端点条件都使用 $o_N$，译文据此改正。}
\[
  L=\phi_H^1(o_N).
\]
\begin{enumerate}[label=(\arabic*)]
  \item 映射
  \[
    \Phi_H:L\cap\nu^*S
    \longrightarrow\operatorname{Chord}(H;o_N,\nu^*S),
  \]
  定义为
  \[
    x\longmapsto z_x^H(t)
    =\phi_H^t(\phi_H^1)^{-1}(x),
  \]
  它在作为常值道路的 $L\cap\nu^*S$ 与 $H$ 的\personHamilton{}方程解集之间给出一一对应。\footnote{译者注：原书把 $\Phi_H$ 的定义域印成 $o_N\cap L$，但同一句随即写明所对应的常值道路集合是 $L\cap\nu^*S$；译文依结论及端点条件改为后者。}

  \item 映射 $a\longmapsto\Phi_H(a)$\footnote{译者注：原书此处的 $a$ 未定义，而第 (1) 项只把 $\Phi_H$ 定义在交点 $x\in L\cap\nu^*S$ 上；要得到两个方程解集之间的对应，还需要给出作用于带映射的坐标变换。原书没有写出该变换，译文不自行重构。}还在方程~\eqref{eq:floer-equation-zero-conormal} 的解集与下列方程的解集之间定义一一对应：
  \begin{equation}
    \begin{cases}
      \dfrac{\partial v}{\partial\tau}
      +J^H\dfrac{\partial v}{\partial t}=0,\\[6pt]
      v(\tau,0)\in\phi_H^1(o_N),
      \qquad v(\tau,1)\in\nu^*S,
    \end{cases}
  \label{eq:dynamical-geometric-solution-correspondence}
  \end{equation}
  其中
  \[
    J^H=\{J_t^H\},
    \qquad
    J_t^H:=\bigl(\phi_H^t(\phi_H^1)^{-1}\bigr)_*J_t.
  \]
  此外，方程~\eqref{eq:dynamical-geometric-solution-correspondence} 正则，当且仅当方程~\eqref{eq:floer-equation-zero-conormal} 正则。
\end{enumerate}
\end{lemma}

把方程~\eqref{eq:floer-equation-zero-conormal} 变换成方程~\eqref{eq:dynamical-geometric-solution-correspondence} 后，还可以把 $J^H$ 进一步形变到常值族 $J_0$，并考虑
\begin{equation}
  \begin{cases}
    \dfrac{\partial v}{\partial\tau}
    +J_0\dfrac{\partial v}{\partial t}=0,\\[6pt]
    v(\tau,0)\in\phi_H^1(o_N),
    \qquad v(\tau,1)\in\nu^*S.
  \end{cases}
\label{eq:geometric-floer-equation-constant-j}
\end{equation}
这个形变保持相伴\personFloer{}复形的滤过；证明可仿照\cref{lem:spectral-number-independent-of-j}的证明。

下述命题给出与方程~\eqref{eq:dynamical-geometric-solution-correspondence}和~\eqref{eq:geometric-floer-equation-constant-j}相伴的作用量泛函；它决定相伴\personFloer{}同调 $HF(L,\nu^*S)$ 的自然滤过。

\begin{proposition}
\label{prop:effective-action-functional}
令 $L=L_H$ 与 $h_H$ 如\cref{lem:basic-generating-function}所述。令 $\Omega(L,\nu^*S;T^*N)$ 是满足
\[
  \gamma(0)\in L,
  \qquad
  \gamma(1)\in\nu^*S
\]
的道路 $\gamma:[0,1]\to T^*N$ 所成的空间。考虑\termEffectiveActionFunctional{}
\[
  \mathcal A^{\mathrm{eff}}(\gamma)
  =\int\gamma^*\theta+h_H(\gamma(0)).
\]
则
\[
  d\mathcal A^{\mathrm{eff}}(\gamma)(\xi)
  =\int_0^1\omega\bigl(\xi(t),\dot\gamma(t)\bigr)\,dt.
\]
特别地，
\begin{equation}
  \mathcal A^{\mathrm{eff}}(c_x)
  =h_H(x)
  =\mathcal A_H^{\mathrm{cl}}(z_x^H)
\label{eq:effective-action-at-constant-path}
\end{equation}
对常值道路 $c_x\equiv x\in L\cap\nu^*S$ 成立，也就是说，对 $\mathcal A^{\mathrm{eff}}$ 的任意临界道路 $c_x$ 成立。\footnote{译者注：原书本命题把 $h_H$ 先写成 $h_L$，把道路值域写成 $\mathbb R$，又把与 $\nu^*S$ 相伴的同调、道路终点和临界交点多次写成 $o_N$（其中还出现“$\gamma(0)\in L,o_N$”）。译文依\cref{lem:basic-generating-function}、道路空间的名称以及方程 (12.7.36)--(12.7.38)，统一为 $h_H$、$T^*N$ 与第二边界 $\nu^*S$。}
\end{proposition}

我们要特别强调，上述有效作用量泛函的定义中出现了“\termBoundaryContribution{}”$h_H(\gamma(0))$。为了使\personCauchy{}--\personRiemann{}方程~\eqref{eq:geometric-floer-equation}或~\eqref{eq:geometric-floer-equation-constant-j}成为相应作用量泛函的\emph{梯度流线方程}，必须加入这一项。

仿照定义 $\rho(H,J;S,a)$ 的方式，可以先对正则的 $\overline J$ 定义有效作用量泛函 $\mathcal A^{\mathrm{eff}}$ 的
\[
  \overline\rho(L,\overline J;S,a),
\]
再证明它与 $\overline J$ 无关。于是由构造可得\footnote{译者注：原书本段先把最后一个参数 $a$ 改成 $u$，又依次使用 $\overline\rho(L,\overline J;u,S)$、$\rho(H,S,J)$ 与 $\overline\rho(L,S,J^H)$；下述定理和命题继续使用缺少上同调类且次序变化的记号。译文统一沿用定义 12.6.5 的 $\rho(H;S,a)$，并相应写成 $\overline\rho(L,\overline J;S,a)$。}
\[
  \rho(H;S,a)=\overline\rho(L,J^H;S,a).
\]

上述讨论立即给出如下推论。

\begin{theorem}
\label{thm:spectral-invariant-depends-on-lagrangian-and-generator}
考虑两个哈密顿量 $H$ 与 $H'$，使得
\[
  \phi_H^1(o_N)=\phi_{H'}^1(o_N),
  \qquad
  h_H=h_{H'},
\]
其中 $h_H$ 与 $h_{H'}$ 分别是 $\phi_H^1(o_N)$ 与 $\phi_{H'}^1(o_N)$ 的相伴生成函数。则
\[
  \rho(H;S,a)=\rho(H';S,a).
\]
\end{theorem}

注意，条件 $\phi_H^1(o_N)=\phi_{H'}^1(o_N)$ 只依赖于哈密顿向量场，或者说哈密顿量的微分；而条件 $h_H=h_{H'}$ 则依赖于哈密顿量本身的实际取值。有一种特别自然的情形，使这两个要求自动成立。

\begin{proposition}
\label{prop:spectral-invariant-time-one-image}
设 $N$ 连通。令 $B\subset N$ 是具有非空内部的子集，并考虑满足
\[
  \operatorname{supp}H,\operatorname{supp}H'
  \subset T^*N\setminus T
\]
的哈密顿量 $H,H'$，其中 $T$ 是 $T^*N$ 中 $o_B$ 的一个\termTubularNeighborhood{}。若
\[
  \phi_H^1(o_N)=\phi_{H'}^1(o_N),
\]
则 $h_H=h_{H'}$。特别地，对这样的哈密顿量 $H$，$\rho(H;S,a)$ 只依赖于时一像 $\phi_H^1(o_N)$。\footnote{译者注：原书只对 $H$ 陈述支集假设，却在条件与证明中使用未量化的 $H'$；证明还把 $o_B\subset T^*N$ 简写成 $B$，写作 $T\supset B$。译文补明 $H'$ 满足同一支集条件，并写回 $T\supset o_B$。}
\end{proposition}

\begin{proof}
由方程
\[
  dh_H=i_H^*\theta
\]
可知，恒等式 $\phi_H^1(o_N)=\phi_{H'}^1(o_N)$ 蕴含
\[
  dh_H=dh_{H'},
\]
因而 $h_H-h_{H'}$ 是常数。由于 $H,H'\equiv0$ 于 $T\supset o_B$，对 $q\in B$ 有
\[
  z_q^H\equiv q\equiv z_q^{H'}.
\]
所以
\[
  h_H(q)=h_{H'}(q)=0.
\]
这就完成了证明。
\end{proof}

\begin{example}
\label{ex:local-hamiltonian-spectral-invariant}
考虑辛流形 $(M,\omega)$，并令 $B\subset M$ 是具有非空内部的子集。设 $H=H(t,x)$ 是 $M$ 上满足
\[
  \operatorname{supp}H\subset M\setminus B
\]
的时变哈密顿量。假设流 $t\longmapsto\phi_H^t$ 在 $C^0$ 意义下很小，因而存在对角线 $\Delta_M\subset M\times M$ 的一个\termDarbouxWeinsteinNeighborhood{} $V_\Delta$，使得
\[
  \operatorname{Graph}\phi_H^t\subset V_\Delta.
\]
若对 $\mathbf x=(x,y)\in V_\Delta$ 定义
\[
  \mathbb H(t,\mathbf x)=H(t,x),
\]
并在 $V_\Delta$ 外借助\termCutOffFunction{}将其延拓到 $T^*\Delta$，那么 $\mathbb H$ 满足\cref{prop:spectral-invariant-time-one-image}所需的全部假设。这样便可以定义\termHamiltonianSpectralInvariant{}的一个局部版本；它在 (Oh11b) 中记为
\[
  \rho_{\mathcal U}^{\mathrm{ham}}(H;a),
\]
并且只依赖于时一映射 $\phi_H^1$ 的图。这里 $\mathcal U\subset\mathcal L_0(M)$ 是由短道路 $\gamma:[0,1]\to M$ 组成的子集，这些道路对所有 $t\in[0,1]$ 都满足\footnote{译者注：原书把单个有序对 $(\gamma(t),\gamma(0))$ 与邻域 $V_\Delta$ 的关系写成子集号 $\subset$；译文改为属于号 $\in$。}
\[
  \bigl(\gamma(t),\gamma(0)\bigr)\in V_\Delta.
\]
\end{example}
